\section{Methodology}
\label{sec:methodology}

\subsection{Teamwork and Collaboration}
\label{subsec:teamwork-collaboration}

Tasks were distributed in pairs throughout most of the project, which supported collaboration and peer review. The workflow was organized in Trello across five completed sprints, keeping implementation, testing, diagram updates, and report work visible as the project evolved. \figref{fig:trello-board} shows the final sprint board used by the group.

\begin{figure}[htbp]
\centering
\optionalgraphic[width=\textwidth]{assets/img/Trello.png}
\caption{Final Trello sprint board used to coordinate implementation, testing, diagrams, and report tasks.}
\label{fig:trello-board}
\end{figure}
\FloatBarrier

During Sprint 1, the work centered on planning the project structure, preparing the project proposal, creating initial diagrams, and designing the database. Sprint 2 moved into setup and early implementation: configuring the application environment, drafting the first user interface, setting up SQL and NoSQL databases, and implementing authentication for patients, doctors, and administrators.

Sprint 3 added audio capture, speech-to-text transcription, the booking and scheduling flow, database fixes, and preparation for the midterm project presentation. Sprint 4 covered document versioning and editing, email sending, LLM-supported transcription, suggestions and report generation, UI completion, server-side PDF generation, logging, system testing, and AOP work. Sprint 5 focused on report completion, booking/admin permission refinements, database and diagram updates, PDF export naming, email attachments for generated PDFs, video recording, LaTeX report rewriting, and supervisor review. \appref{app:meeting-log} gives the detailed meeting log behind this sprint summary and maps each recorded meeting to the relevant sprint.

GitHub was used to share and manage the code in one shared repository\footnote{\url{\RepositoryUrl}}, keeping contributions organized while supporting collaboration. The group met around twice per week to review progress, coordinate tasks, and resolve blockers. WhatsApp handled daily updates, while Discord was used for sharing files, data, and other project resources.

\subsection{Scope Context and Source Selection}
\label{subsec:scope-context-source-selection}

During early scoping, the project was narrowed around clinical documentation rather than general medical administration, image analysis, or unrelated laboratory AI. This direction is supported by studies showing that physicians spend a large amount of time on electronic records and clerical work during and after consultations\refcite{1,2,3}. The team therefore selected a problem where software engineering could address a concrete workflow: capturing consultation audio, converting it to text, generating a draft note, and supporting a structured review step.

The Danish context also shaped the scope. The Danish digital health strategy describes a healthcare system that should feel coherent and trustworthy for patients while making daily workflows easier for healthcare professionals\refcite{7}. At SDU and Odense University Hospital, CAI-X presents clinical AI as a collaboration between healthcare staff, researchers, and engineers, with solutions chosen from the needs and challenges clinicians face in daily work\refcite{8}. This gave the project a useful local reference point for treating AI as workflow support instead of as a replacement for clinical judgement.

The most directly related Danish example found during source selection is CAI-X's AID\_NOTE project\refcite{9}, which evaluates Ambient Scribe speech-to-text solutions across clinical settings and measures documentation quality, time consumption, user experience, and economic impact. Med Flow is not a deployment of that project, but it addresses a comparable type of problem at semester-project scale. The technical source selection therefore focused on risk guidance for AI medical scribes\refcite{4}, robust local speech recognition\refcite{10}, and responsible AI governance\refcite{11,12}.

\subsection{Application of Course Knowledge}
\label{subsec:course-knowledge}

\begin{itemize}
	\item \textbf{Artificial Intelligence:} Faster-Whisper converts consultation audio into text, while the local Qwen3:8b model generates structured report drafts, prescription suggestions, and Suggestive Mode alerts. The AI work is tied to explicit safeguards: prompt grounding, JSON-schema validation, hallucination checks, ASR artifact filtering, and doctor approval before clinical output becomes final.
	\item \textbf{Component-Based Systems:} FastAPI routers, application services, domain contracts, repositories, and infrastructure adapters separate authentication, patients, appointments, consultations, transcription, review, prescriptions, PDF, email, and audit. This makes it possible to replace Whisper, Ollama, persistence adapters, email delivery, or booking integration without rewriting the whole workflow.
	\item \textbf{Operating Systems:} Docker Compose was used to containerize the FastAPI application, Whisper sidecar, MySQL, MongoDB, MailHog, and supporting services. This connected the operating-systems course knowledge to the project through process isolation, port mapping, filesystem volumes, environment variables, and reproducible local deployment across different developer machines.
	\item \textbf{Algorithms and Data Structures:} the system uses status transitions for consultations and appointments, transcript chunk assembly, patient and appointment filtering, version-history lists for edited documents, deterministic safety rules for prescriptions, and structured JSON objects validated before they enter the review flow. These structures turn raw speech and model output into data that can be searched, reviewed, tested, and rendered as PDFs.
	\item \textbf{Designing Software Systems with Embedded Elements:} although the application is not deployed on a microcontroller, it follows embedded-oriented constraints by running AI services locally, budgeting memory and VRAM, validating CPU/GPU fallback paths, containerizing services with Docker Compose, and avoiding cloud APIs for sensitive processing. The hardware envelope shows that the design is constrained by real local machines rather than unlimited server resources.
	\item \textbf{Data Management:} MySQL stores users, patients, appointments, consultations, prescriptions, and audit logs as structured records, while MongoDB stores transcripts, generated documents, prompts, email templates, and AI artifacts. The split supports relational consistency for approved clinical data and schema flexibility for iterative AI outputs.
\end{itemize}